\documentclass[11pt]{article}
\usepackage[margin=0.75in]{geometry}
\usepackage[T1]{fontenc}
\usepackage[utf8]{inputenc}
\usepackage{lmodern}
\usepackage{microtype}
\usepackage{array}
\usepackage{booktabs}
\usepackage{longtable}
\usepackage{tabularx}
\usepackage{enumitem}
\usepackage{needspace}
\usepackage{multicol}
\usepackage{xcolor}
\usepackage{hyperref}
\usepackage[most]{tcolorbox}
\usepackage{tikz}
\usetikzlibrary{arrows.meta,calc,positioning,shapes.geometric}

% Semantic color names are retained for compatibility with the weekly files,
% but every value is pure black or white for monochrome printing.
\definecolor{rubricred}{gray}{0}
\definecolor{linkblue}{gray}{0}
\definecolor{softblue}{gray}{1}
\definecolor{softgold}{gray}{1}
\definecolor{softgray}{gray}{1}

\hypersetup{
  hidelinks,
  pdfborder={0 0 0}
}

\setlength{\parindent}{0pt}
\setlength{\parskip}{0.55em}
\setlength{\tabcolsep}{3pt}
\setlist[itemize]{leftmargin=1.25em, itemsep=0.2em}
\renewcommand{\arraystretch}{1.22}

\newcommand{\weektitle}[3]{
  \begin{center}
    {\Large\bfseries #1}\\[0.2em]
    {\large #2}\\[0.2em]
    {\small #3}
  \end{center}
  \vspace{0.5em}
}

\newcommand{\properrow}[4]{\textbf{#1} & #2 & #3 & #4\\}
\newcommand{\properrefs}[1]{\textnormal{(\textit{#1})}}
\newcommand{\sourceurl}[2]{\href{#1}{#2}}
\newcommand{\latin}[1]{\textit{#1}}
\newcommand{\rubric}[1]{\textbf{#1}}

\newenvironment{properstable}
{\footnotesize\begin{longtable}{>{\raggedright\arraybackslash}p{0.15\linewidth}>{\raggedright\arraybackslash}p{0.25\linewidth}>{\raggedright\arraybackslash}p{0.18\linewidth}>{\raggedright\arraybackslash}p{0.34\linewidth}}
\textbf{Proper} & \textbf{TLM text / reference} & \textbf{Scriptural axis} & \textbf{Connection}\\
\midrule
\endhead}
{\bottomrule\end{longtable}}

\newenvironment{sensestable}
{\vspace{-0.35em}\begingroup\footnotesize\renewcommand{\arraystretch}{1.10}%
\noindent\hspace*{\dimexpr0.04\linewidth-4\tabcolsep\relax}%
\begin{tabular}{>{\bfseries\raggedright\arraybackslash}p{0.15\linewidth}>{\raggedright\arraybackslash}p{\dimexpr0.77\linewidth+4\tabcolsep\relax}}
\textbf{Sense} & \textbf{Synthesis}\\
\midrule}
{\bottomrule\end{tabular}\par\endgroup}

\newenvironment{connectionstable}
{\small\begin{longtable}{>{\bfseries\raggedright\arraybackslash}p{0.18\linewidth}>{\raggedright\arraybackslash}p{0.74\linewidth}}
\toprule
\textbf{Thread} & \textbf{Synthesis}\\
\midrule
\endhead}
{\bottomrule\end{longtable}}

\newenvironment{witnessstable}
{\small\begin{longtable}{>{\raggedright\arraybackslash}p{0.2\linewidth}>{\raggedright\arraybackslash}p{0.34\linewidth}>{\raggedright\arraybackslash}p{0.38\linewidth}}
\toprule
\textbf{Witness} & \textbf{Text / source status} & \textbf{Interpretive contribution}\\
\midrule
\endhead}
{\bottomrule\end{longtable}}

\newcommand{\witnessrow}[3]{#1 & #2 & #3\\}
\newcommand{\sectionguard}{\Needspace{7\baselineskip}}

% Keep the opening survey page compact but unchanged; sustained analysis uses
% a slightly smaller reading size so longer patristic dossiers remain usable.
\newenvironment{deepstudy}
  {\begingroup\small}
  {\par\endgroup}

\newenvironment{threestable}[3]
{\small\begin{longtable}{>{\bfseries\raggedright\arraybackslash}p{0.18\linewidth}>{\raggedright\arraybackslash}p{0.34\linewidth}>{\raggedright\arraybackslash}p{0.40\linewidth}}
\toprule
\textbf{#1} & \textbf{#2} & \textbf{#3}\\
\midrule
\endhead}
{\bottomrule\end{longtable}}

\newenvironment{fourstable}[4]
{\footnotesize\begin{longtable}{>{\bfseries\raggedright\arraybackslash}p{0.15\linewidth}>{\raggedright\arraybackslash}p{0.23\linewidth}>{\raggedright\arraybackslash}p{0.25\linewidth}>{\raggedright\arraybackslash}p{0.27\linewidth}}
\toprule
\textbf{#1} & \textbf{#2} & \textbf{#3} & \textbf{#4}\\
\midrule
\endhead}
{\bottomrule\end{longtable}}

\newtcolorbox{studybox}[1]{
  enhanced,
  breakable,
  colback=white,
  colbacktitle=white,
  coltitle=black,
  frame hidden,
  boxrule=0pt,
  sharp corners,
  borderline west={1.2pt}{0pt}{black},
  left=8pt,
  right=2pt,
  top=2pt,
  bottom=2pt,
  toptitle=0pt,
  bottomtitle=2pt,
  before skip=8pt,
  after skip=8pt,
  title={#1},
  fonttitle=\small\bfseries
}

\newtcolorbox{massbox}[1]{
  enhanced,
  breakable,
  colback=white,
  colframe=black,
  colbacktitle=white,
  coltitle=black,
  boxrule=0.55pt,
  sharp corners,
  left=7pt,
  right=7pt,
  top=4pt,
  bottom=4pt,
  toptitle=3pt,
  bottomtitle=3pt,
  titlerule=0.45pt,
  before skip=8pt,
  after skip=8pt,
  title={#1},
  fonttitle=\small\bfseries
}

\newtcolorbox{sourcecard}[1]{
  enhanced,
  colback=white,
  colframe=black,
  colbacktitle=white,
  coltitle=black,
  boxrule=0.4pt,
  sharp corners,
  left=5pt,
  right=5pt,
  top=4pt,
  bottom=4pt,
  toptitle=3pt,
  bottomtitle=3pt,
  titlerule=0.4pt,
  before skip=4pt,
  after skip=4pt,
  title={#1},
  fonttitle=\small\bfseries
}

\tikzset{
  guide node/.style={draw=black, fill=white, align=center, inner sep=5pt, font=\small},
  guide emphasis/.style={draw=black, very thick, fill=white, align=center, inner sep=5pt, font=\small},
  guide arrow/.style={-{Stealth[length=2mm]}, thick, draw=black}
}

\hypersetup{pdftitle={The Four Last Things: The Truth of a Life},pdfsubject={Judgment: a Catholic study of particular and final judgment, grace, works, mercy, and justice}}
\setcounter{tocdepth}{1}\setlength{\emergencystretch}{2em}
\input{articles/faith/four-last-things-shared/format}
\begin{document}
\ifdefined\TriptychPrintEdition
\FourLastThingsTitle{The Truth of a Life}{Judgment}
\else
\fi
\FourLastThingsNavigation
\section{The Court We Already Desire}

Human beings often protest judgment until injustice reaches them. Then neutrality appears not merciful but monstrous. The betrayed want the betrayal named; the forgotten dead demand that history be more than the victor's account. Judgment promises that wrong will finally be answered. It also threatens the disguises by which we excuse ourselves.

Christian judgment is not the universe's final administrative procedure. The judge is Jesus Christ, crucified and risen (Jn 5:22--29; Acts 17:31). The wounds of the Judge prevent two evasions at once: justice is not indifferent to what sin costs, and mercy is not the indulgence of one untouched by the cost. The one before whom every life is disclosed is the one who bore sin and offered reconciliation.

\section{Particular and Final}

Catholic teaching distinguishes the particular judgment from the last judgment. At death each person receives eternal retribution in the immortal soul as his life is referred to Christ: heaven immediately or through purification, or definitive loss (CCC 1021--1022). At Christ's return the dead rise, history reaches its consummation, and the truth of hearts and deeds is manifested in their relations and consequences (Mt 25:31--46; CCC 1038--1041).

The person does not disappear between these judgments. Consciousness and will subsist in the surviving soul; the human self continues while awaiting bodily resurrection (CDF, letter on eschatology, 1979, nos.~2--5). The same person whose destiny is settled at death will stand risen before Christ. The last judgment brings no second opportunity to reverse a definitive refusal of grace and no new defendant in place of the old.

What, then, does it add? First, bodily consummation: those who lived and acted through bodies receive their destiny as whole, risen human beings. Second, public truth: each life is manifested within the completed history it helped to shape. Particular judgment does not exhaust either the body's destiny or the revelation of God's providence. The saints already enjoy God, yet they still await resurrection and the gathering of the whole redeemed creation.

\section{The Judge Is the One Judged by Us}

The Father has given judgment to the Son; the appointed judge is the man God raised from the dead (Jn 5:22--29; Acts 17:31). Human life is measured against the truth of humanity made visible in the obedient Son. His authority is inseparable from his saving work: the one who judges has first called sinners into reconciliation with God.

That judge was himself arraigned, falsely witnessed against, mocked, and condemned. This does not make the last judgment a reversal in which Jesus finally enjoys the power his enemies used. It reveals the difference between divine and fallen judgment. Pilate protects order by surrendering the innocent; Caiaphas turns political calculation into necessity; the crowd converts unanimity into truth. Christ judges without needing a victim to stabilize his authority. His wounds disclose both the world's verdict upon God and God's mercy toward the world.

Matthew 25 therefore refuses a merely interior religion. The king identifies himself with those whose hunger, estrangement, nakedness, illness, and imprisonment exposed them to neglect. The passage does not say that every policy dispute has an obvious answer, nor that works of mercy save apart from Christ. It says something more searching: relation to Christ becomes visible in relation to persons who cannot increase our prestige.

\section{Faith, Works, and Grace}

Scripture speaks of judgment according to works (Rom 2:6--11; 2~Cor 5:10; Rev 20:12) and denies that anyone saves himself by works performed apart from grace. Ephesians 2:8--10 holds the two truths together within a single movement: salvation is God's gift, and the saved are created in Christ for good works. Grace does not merely excuse a life from being assessed. It heals freedom so that a different life becomes possible.

The sheep in Matthew 25 do not purchase heaven by philanthropic arithmetic. Their treatment of the hungry, stranger, sick, and imprisoned manifests whether Christ's love has taken flesh in them. James's claim that faith without works is dead (Jas 2:14--26) does not oppose faith to grace. It opposes living faith to verbal possession.

The Catechism makes the order precise. No one merits the initial gift of forgiveness and justification. Once adopted into Christ, a person can act freely under the Spirit's initiative, and these acts of charity have genuine merit before God (CCC 2008--2011). Good works are therefore more than outward evidence that something has already happened inside us: they belong to the grace-enabled growth of love and its promised reward. God crowns gifts that have truly become his children's acts. Judgment honors both the priority of grace and the reality of the response it makes possible.

\section{What Is Disclosed, and Why Publicly?}

The objection to a second judgment is serious: if God already knows everything and the soul already receives its eternal destiny, why await another judgment? Aquinas puts that objection directly in ST III, q.~59, a.~5. His answer distinguishes the end of a person's earthly choosing from the unfinished history to which that life belongs. Reputation can remain false after death; teaching can continue to help or corrupt; the bodies of the dead can be honored or abused. A person's life has ended, but everything concerning that life has not yet reached its conclusion.

In his reply to the first objection, Aquinas expressly denies that the dead acquire new merit or demerit through these continuing effects. His point is that the whole truth of their lives belongs under God's judgment, including relations and consequences only completed at history's end. A child is not condemned for a parent's choices, nor does a later reader's abuse of a book automatically become its dead author's new sin. God judges each person's own freedom within the real history of its exercise. Aquinas supplies a theological explanation of the Church's teaching that final judgment reveals deeds to their furthest consequences and makes God's providence known (CCC 1039--1040).

This public truth matters especially where earthly judgment has failed. A victim whose testimony was suppressed will not remain dependent on the oppressor's account; unnoticed fidelity will not disappear because no archive recorded it. Yet disclosure must include mercy as truthfully as guilt. Thomas holds that even forgiven sins are disclosed, manifesting mercy rather than bringing the blessed shame (Supplement, q.~87, a.~2). He presents this as the more probable opinion, noting Peter Lombard's contrary view. The Church promises the manifestation of divine justice, not a spectacle modeled on human humiliation. We are not given the event's perceptual mechanism or a warrant to imagine the blessed endlessly reliving shame.

\section{The Unequal Histories of Freedom}

If judgment is according to deeds, does it ignore the unequal conditions under which people act? One person receives sound formation, stable affection, and ready access to the sacraments; another inherits violence, deception, compulsion, or ignorance. Catholic moral theology does not solve this by declaring acts meaningless. It distinguishes the objective character of an act from the agent's culpability, which can be diminished by ignorance, fear, habit, duress, and other factors (CCC 1735, 1750--1756, 1860).

Only God knows these histories without simplification. This is not a loophole by which every grave act becomes inculpable, nor a warrant to suspend necessary human justice. Courts and communities must judge acts within their competence, protect victims, and seek truth by proportionate evidence. They may not pronounce the eternal measure of a person. Final judgment consoles precisely because neither privilege nor trauma is invisible to it.

Grace also prevents unequal histories from becoming fate. No one converts by unaided power, and no wound places a person outside God's possible action. Yet grace is not coercion wearing a holy name. It heals and elevates agency so that consent is genuinely the person's. The Catholic refusal of both self-salvation and deterministic theater protects the drama judgment reveals: God acts first and throughout; the creature's graced answer is real.

\section{Conscience Without Self-Acquittal}

Conscience is not the final court. It is a judgment of practical reason that must seek truth and may err (CCC 1776--1794). A quiet conscience can result from innocence, but also from ignorance cultivated for convenience. A tormented conscience can identify guilt, but also suffer scrupulosity. ``I feel no guilt'' and ``I feel unforgivable'' are both insufficient verdicts.

The prospect of judgment should lead to concrete examination, repentance, reparation, and trust in received mercy. Presumption refuses the need for conversion. Scrupulous distress can instead keep reopening what has been honestly confessed; such distress should not be treated as proof of guilt or a voluntary refusal to trust. It calls for patient pastoral care, not an added accusation. The soul's Judge is Christ; intensity of fear supplies no privileged knowledge of his verdict.

\section{Mercy at the Root of Justice}

Divine mercy does not persuade God to cease being just. Mercy confronts misery and acts to overcome it; justice gives what is due within the order God wisely creates. Aquinas asks why both are present in God's works (ST I, q.~21, a.~4). Whatever is due to a creature presupposes a gift already made: its existence, its nature, or its calling. Traced to their source, even the claims justice honors depend on God's freely given goodness. Mercy is therefore the root of divine justice, not an exception purchased by persuading God to overlook the truth. At the Cross, forgiveness condemns sin in the flesh and opens the sinner's return (Rom 8:3).

Universal salvific will and real judgment must both be confessed. God desires all to be saved (1~Tim 2:4); Christ's sacrifice is sufficient for all; grace precedes and enables every saving response. Yet Scripture's summons to repent is not theater, and freedom is not redeemed by being abolished. The Church prays for the salvation of all; these teachings give us neither a guarantee that all will be saved nor authority to pronounce on another person's damnation.

\section{Purgatory and the Refusal of Cheap Acquittal}

The final alternatives are communion with God and definitive loss. Purgatory belongs wholly on the side of salvation: those who die in God's grace and friendship but remain imperfectly purified are assured of heaven (CCC 1030--1032). Forgiveness can be real while attachment to sin and its disorder still require healing. The purified are not being invited to reverse a final refusal of God; they are being made wholly ready for the friendship they have received. Christ's victory accomplishes this healing.

Benedict XVI considers the proposal that the purifying fire is the encounter with Christ himself (\emph{Spe salvi} 47). He develops it as a way to understand transformation: the Judge's love exposes and heals what cannot remain in communion with God. The passage offers a theological account without defining a mechanism or an earthly duration. Its force is that mercy makes us truthful and whole; it does not simply attach a favorable verdict to an unchanged person.

But if Christ purifies and the person's destiny is settled, what can our prayer contribute? Benedict takes up precisely that question in the next paragraph (\emph{Spe salvi} 48): persons belong to one another, and intercession is a real expression of their communion in Christ. The Church teaches that prayer, especially the Eucharistic sacrifice, assists the departed in their purification (CCC 1032). These suffrages are help offered within salvation, not a new conversion imposed on someone who has definitively refused God. Their efficacy comes from Christ, who includes the love of his members in his care for them. All Souls thus gives the living a work of charity: to commend the dead, have Mass offered for them, and entrust the fruit to God.

\section{Living Under the True Verdict}

The Christian prepares for judgment by refusing both theatrical terror and practical forgetfulness. The decisive habits are plain: remain in grace, confess grave sin, forgive, restore what injustice took, feed the person in reach, pray, keep promises, and ask for the grace of perseverance. Such acts do not manipulate the Judge. They conform life to him.

The final word is not our self-description. It is Christ's truthful judgment. That fact wounds pride and frees the conscience from endless self-invention. The question is not whether we can construct an acquittal. It is whether we will come into the light while mercy is offered (Jn 3:19--21).

The liturgy makes this future judgment present without pretending to stage it. The Creed confesses that Christ will come to judge the living and the dead; the Eucharistic liturgy remembers his death, proclaims his Resurrection, and awaits his coming. Worship trains judgment's proper posture: confession before self-defense, mercy received before mercy given, and hope fixed on the returning Christ.

\appendix
\section{Scope, Method, and Qualifications}
This article distinguishes particular judgment from the embodied and public consummation of final judgment. It treats works, grace, merit, conscience, mercy, purification, and suffrages for the dead. Scripture and authoritative Catholic teaching govern; Aquinas's explanations and Benedict XVI's account of purifying encounter retain their identified theological status. Pastoral applications are synthesis. The article neither calculates personal culpability nor adjudicates the Catholic schools on grace and predestination, and its brief treatment of scrupulosity does not replace individual care. No person's eternal state or independent dogmatic, Thomistic, pastoral, psychological, or ecclesiastical review is claimed.
\input{articles/faith/four-last-things-02-judgment/references}
{\scriptsize\input{articles/faith/four-last-things-02-judgment/generation-metadata}}
\end{document}
